% Generado por sistema/backend ReporteTex.java -- no editar a mano.
% Datos: /private/tmp/claude-501/-Users-miguel-project-taller-s/a868368b-4351-4d87-8017-fb9a5bc093aa/scratchpad/titulados.csv (titulados) / /private/tmp/claude-501/-Users-miguel-project-taller-s/a868368b-4351-4d87-8017-fb9a5bc093aa/scratchpad/empleadores.csv (empleadores)
\begin{enumerate}
  \item \textbf{Condiciones de empleo (3.2.3).} 3 de 3 indicadores del criterio son favorables: ``Titulados con empleo o emprendimiento'' (83.1\%); ``Consideran pertinentes a su formación las labores de su puesto'' (90.7\%); ``Perciben concordancia entre su formación y lo que pide el mercado'' (61.9\%). El tiempo medio para conseguir el primer empleo es de 2.63 meses (IC 95\%: 1.85 a 3.42), y queda como línea base porque ARCU-SUR no fija un valor de referencia. El 80.6\% de quienes trabajan en una organización lo hace en el sector privado, principalmente en ``Desarrollo de Software''. La mediana de experiencia es de 6 años (media 8.95); ocupa un cargo ejecutivo el 22.6\% de los consolidados que declararon un cargo, frente al 7.5\% de los junior.
  \item \textbf{Resultados de la formación (3.2.1) y actualización curricular (2.1.7).} El 74.1\% está satisfecho con la formación recibida, pero el 71.6\% señala ``Actualizar el plan de estudios'' como el aspecto a mejorar más importante. ``Base de Datos'' es la asignatura que más ventaja da en el trabajo (94.6\%), y los empleadores piden reforzar sobre todo ``Inteligencia artificial y aprendizaje automático'' y ``Ciberseguridad y seguridad de la información''.
  \item \textbf{Perfil de egreso (2.1.2).} 7 de 7 ítems del perfil son favorables: ``La Carrera brinda confianza como formadora'' (88.7\%); ``El título es consistente con las competencias de los titulados'' (86.8\%); ``Los empleadores conocen el perfil profesional declarado'' (84.9\%); ``El perfil de egreso es coherente con el desempeño de los titulados'' (79.2\%); ``La formación permite a los titulados un desempeño destacado'' (88.7\%); ``Los titulados muestran competencias laborales'' (86.8\%); ``Los titulados muestran valores y ética profesional'' (83.0\%).
  \item \textbf{Posgrado (2.4.1 y 1.1.5).} El 83.1\% quiere un posgrado, sobre todo en ``IA'' (43.9\% de los interesados); las dos áreas más elegidas reúnen el 67.5\% de las preferencias, frente al 70\% que el enunciado toma como referencia para abrir programas modulares. El 70.9\% considera afín la oferta de la FCyT y el 41.5\% de los interesados la elegiría; aspectos desfavorables del criterio: ``Interesados que elegirían la UMSS'' (41.5\%). El 91.8\% de quienes ya cursaron un programa lo pagó con recursos propios.
  \item \textbf{Relación con las organizaciones (2.4.2).} 6 de 12 aspectos son favorables: ``Organizaciones que contrataron titulados en los últimos 5 años'' (79.2\%); ``La Carrera considera a los empleadores en prácticas y pasantías'' (58.5\%); ``Los empleadores participaron en procesos de retroalimentación'' (52.8\%); ``La Carrera mantiene vínculos que generan oportunidades'' (62.3\%); ``Los empleadores participaron en vinculación entre profesionales'' (50.9\%); ``Las organizaciones recurren a la UMSS para contratar'' (60.4\%). Con menos de la mitad de acuerdo: ``Las autoridades de la Carrera consultan a los empleadores'' (47.2\%); ``Los empleadores participaron en vinculación con el entorno laboral'' (47.2\%); ``Los empleadores participaron en difusión de oportunidades laborales'' (47.2\%); ``La Carrera consultó a los empleadores sobre el desempeño de los titulados'' (47.2\%); ``La Carrera consultó a los empleadores sobre la pertinencia de la formación'' (43.4\%); ``La Carrera consultó a los empleadores en su autoevaluación'' (35.8\%).
  \item \textbf{Seguimiento a graduados (3.2.2).} El 57.4\% quiere formar parte de una red de contactos de la Carrera, y la herramienta permite recibir respuestas de forma continua y recalcular los indicadores.
  \item \textbf{Cruces de variables.} 6 de los 14 cruces muestran dependencia significativa al 5\% (se rechaza $H_0$):
  \begin{itemize}[nosep]
    \item \emph{Nivel del programa cursado $\times$ fuente de financiamiento (C7):} ``Financiamiento'' depende de ``Nivel cursado''; el porcentaje que respondió ``Recursos propios'' en cada grupo de ``Nivel cursado'' es: Diplomado o Especialidad 95.7\%; Maestría o superior 82.8\% (1 de 4 celdas con $FE<5$, conclusión tentativa).
    \item \emph{Edad $\times$ área de interés del posgrado (C8):} ``Área de interés'' depende de ``Edad''; el porcentaje que respondió ``Software y Base de Datos'' en cada grupo de ``Edad'' es: 26 o menos 52.9\%; 27–34 21.1\%; 35 o más 25.5\%.
    \item \emph{Experiencia profesional $\times$ remuneración (C9):} ``Remuneración'' depende de ``Experiencia''; el porcentaje que respondió ``Hasta Bs 5.000'' en cada grupo de ``Experiencia'' es: Junior 56.6\%; Consolidado 5.7\%.
    \item \emph{Experiencia profesional $\times$ cargo (C10):} ``Cargo'' depende de ``Experiencia''; el porcentaje que respondió ``Cargo ejecutivo'' en cada grupo de ``Experiencia'' es: Junior 7.5\%; Consolidado 22.6\%.
    \item \emph{Formación complementaria $\times$ interés en posgrado (C11):} ``Interés en posgrado'' depende de ``Formación complementaria''; el porcentaje que respondió ``No'' en cada grupo de ``Formación complementaria'' es: Sí 12.2\%; No 26.0\%.
    \item \emph{Edad $\times$ situación laboral (C12):} ``Situación laboral'' depende de ``Edad''; el porcentaje que respondió ``No trabaja'' en cada grupo de ``Edad'' es: 26 o menos 41.0\%; 27–34 8.9\%; 35 o más 7.8\% (2 de 9 celdas con $FE<5$, conclusión tentativa).
  \end{itemize}
  En los otros 8 no se rechaza la independencia: Situación laboral $\times$ satisfacción con la formación (C1), Sector $\times$ pertinencia del puesto (C2), Formación complementaria $\times$ remuneración (C3), Interés en posgrado $\times$ oferta de posgrado FCyT (C4), Área de interés $\times$ institución elegida (C5), Concordancia formación–mercado $\times$ tiempo al primer empleo (C6), Género $\times$ interés en posgrado (C13), Formación complementaria $\times$ situación laboral (C14).
\end{enumerate}
